\apendice{Documentación técnica de programación}

\section{Introducción}

El presente anexo tiene como objetivo proporcionar la información técnica necesaria para comprender la estructura interna del sistema desarrollado y facilitar futuras tareas de mantenimiento, ampliación o adaptación del proyecto.

En este documento se describen los principales componentes software que forman parte de la solución implementada, la organización de los directorios del proyecto, los procedimientos de instalación y ejecución, así como las pruebas realizadas para verificar el correcto funcionamiento de los diferentes módulos.

La documentación técnica incluida permite comprender la arquitectura interna del sistema SIEM desarrollado, así como las relaciones existentes entre los diferentes componentes que intervienen en el procesamiento, almacenamiento y visualización de la información.

Asimismo, se proporciona información orientada a futuros desarrolladores que deseen continuar la evolución del proyecto, incorporar nuevas funcionalidades o adaptar el sistema a otros entornos de monitorización. La arquitectura modular utilizada facilita estas tareas y permite realizar modificaciones sobre componentes concretos sin afectar significativamente al funcionamiento global de la aplicación.

\section{Estructura de directorios}

El proyecto se encuentra organizado siguiendo una estructura modular que permite separar claramente las diferentes responsabilidades del sistema.

Esta organización facilita el mantenimiento del código, la incorporación de nuevas funcionalidades y la reutilización de componentes sin afectar al funcionamiento global de la aplicación.

La estructura principal del proyecto es la siguiente:
\begin{verbatim}
	TFG_SIEM/
	|
	+-- api/
	+-- config/
	+-- correlation/
	+-- dashboard/
	+-- db/
	+-- docs/
	+-- ingestor/
	+-- logs/
	+-- machine/
	+-- notifications/
	+-- sensors/
	|
	+-- main.py
	+-- README.md
	+-- requirements.txt
	+-- start.bat
\end{verbatim}

La Figura D.1 muestra la estructura general del proyecto visualizada desde el entorno de desarrollo Visual Studio Code.

\begin{figure}[H]
	\centering
	\includegraphics[width=0.35\textwidth]{vscode_estructura}
	\caption{Estructura general del proyecto visualizada desde Visual Studio Code.}
	\label{fig:vscode_estructura}
\end{figure}

\begin{itemize}
	\item \textbf{api}: Contiene la implementación de la API REST desarrollada con FastAPI.
	\item \textbf{config}: Incluye los parámetros de configuración utilizados por el sistema, tales como umbrales ambientales, rutas de almacenamiento y tiempos de supervisión.
	\item \textbf{correlation}: Implementa el motor de correlación y las reglas encargadas de generar alertas a partir de los eventos procesados.
	\item \textbf{dashboard}: Contiene la aplicación web desarrollada con Streamlit para la visualización de eventos, alertas y anomalías.
	\item \textbf{db}: Gestiona el acceso a la base de datos SQLite y el almacenamiento persistente de eventos y alertas.
	\item \textbf{docs}: Contiene la documentación generada durante el desarrollo del proyecto, incluyendo la memoria y documentación auxiliar.
	\item \textbf{ingestor}: Implementa el módulo encargado de leer y normalizar los logs generados por los dispositivos simulados.
	\item \textbf{logs}: Almacena los archivos de eventos utilizados durante la simulación del entorno monitorizado.
	\item \textbf{machine}: Incluye los mecanismos de detección de anomalías y cálculo de niveles de riesgo.
	\item \textbf{notifications}: Gestiona la simulación de notificaciones asociadas a alertas críticas.
	\item \textbf{sensors}: Implementa la supervisión del estado de los sensores y la detección de dispositivos inactivos o que han dejado de transmitir información.
	\item \textbf{main.py}: Punto de entrada principal de la aplicación. Inicializa los módulos del sistema, ejecuta el proceso de ingesta, lanza las reglas de correlación y coordina el funcionamiento general del SIEM.
	\item \textbf{README.md}: Documento principal de presentación del proyecto donde se describe su funcionamiento, requisitos e instrucciones básicas de instalación y uso.
	\item \textbf{requirements.txt}: Contiene las dependencias necesarias para la instalación y ejecución del proyecto.
	\item \textbf{start.bat}: Script de ejecución utilizado para iniciar automáticamente los principales componentes del sistema en entornos Windows.
\end{itemize}

La separación de responsabilidades entre módulos permite mantener una arquitectura desacoplada y facilita tanto las tareas de mantenimiento como la incorporación de nuevas funcionalidades. Este diseño modular ha permitido desarrollar y validar cada componente de forma independiente antes de su integración dentro del sistema completo.

\section{Manual del programador}

El desarrollo del proyecto se ha realizado utilizando Python como lenguaje principal de programación. La aplicación se ha diseñado siguiendo una arquitectura modular con el objetivo de facilitar su mantenimiento, escalabilidad y reutilización de componentes.

Cada uno de los módulos implementados dispone de una responsabilidad claramente definida, permitiendo realizar modificaciones o ampliaciones sin afectar significativamente al resto del sistema.

El núcleo de la aplicación se encuentra en el fichero \textit{main.py}, encargado de coordinar la lectura de eventos, el almacenamiento de información, la ejecución del motor de correlación y la detección de anomalías.

Para incorporar nuevas reglas de correlación, el desarrollador debe actuar sobre los módulos contenidos dentro del directorio \textit{correlation}. En este componente se implementan las condiciones encargadas de generar alertas a partir de los eventos procesados por el sistema.

Las nuevas reglas de detección de anomalías pueden añadirse dentro del directorio \textit{machine}, donde se encuentran los mecanismos encargados de analizar comportamientos potencialmente anómalos y asignar los niveles de riesgo correspondientes.

La API REST puede ampliarse incorporando nuevos endpoints dentro del directorio \textit{api}. Gracias a la utilización de FastAPI, la incorporación de nuevos servicios resulta sencilla y permite ampliar las capacidades de integración con aplicaciones externas.

Las funcionalidades relacionadas con la visualización de información se encuentran implementadas en el directorio \textit{dashboard}. Desde este módulo es posible añadir nuevas páginas, gráficos o elementos de monitorización utilizando las capacidades proporcionadas por Streamlit.

Los parámetros generales de funcionamiento se almacenan en el directorio \textit{config}, permitiendo modificar umbrales ambientales, tiempos de supervisión o configuraciones generales sin necesidad de alterar el resto del código fuente.

Esta organización facilita el mantenimiento del proyecto y simplifica la incorporación de futuras mejoras o nuevas funcionalidades.

\subsection{Módulo de ingesta}

El módulo de ingesta es el encargado de procesar los archivos de log generados por los diferentes dispositivos simulados. Su función principal consiste en leer los eventos registrados, normalizar su formato y preparar la información para su posterior almacenamiento y análisis.

La implementación permite realizar lecturas incrementales de los archivos de log, evitando reprocesar eventos ya tratados anteriormente.

La Figura D.2 muestra una ejecución real del módulo de ingesta durante las pruebas realizadas, donde puede observarse la lectura de eventos, la normalización de la información y la detección de líneas con formato incorrecto.

\imagen{ingestor}
{Ejecución del módulo de ingesta durante las pruebas realizadas.}

\subsection{Motor de correlación}

El motor de correlación constituye uno de los componentes principales del sistema. Este módulo analiza los eventos normalizados por el ingestor y aplica diferentes reglas de seguridad para identificar situaciones potencialmente peligrosas.

Las reglas implementadas permiten generar alertas ante aperturas de puertas, detección de movimiento, valores ambientales fuera de rango, fallos de comunicación en dispositivos y fallos críticos del sistema eléctrico.

Cada alerta generada incluye información relevante para su posterior análisis, como el tipo de alerta, severidad, nivel de riesgo, estado, depósito afectado, dispositivo implicado y mensaje descriptivo.

\imagen{codigo_correlacion}
{Implementación de las reglas principales del motor de correlación.}

\subsection{Gestión de base de datos}

La persistencia de la información se realiza mediante SQLite. El componente \textit{DatabaseManager} centraliza todas las operaciones relacionadas con el almacenamiento y recuperación de eventos y alertas.

Esta capa de acceso permite desacoplar la lógica de negocio de la gestión de datos, facilitando futuras migraciones a otros motores de bases de datos.

La base de datos implementada en el proyecto se compone principalmente de las tablas \textit{events} y \textit{alerts}, encargadas de almacenar los eventos procesados por el sistema y las alertas generadas por el motor de correlación.

Las Figuras D.4 y D.5 muestran el contenido de ambas tablas durante las pruebas realizadas.

\imagen{sqlite_events}
{Contenido de la tabla \textit{events} almacenada en SQLite.}

\imagen{sqlite_alerts}
{Contenido de la tabla \textit{alerts} almacenada en SQLite.}

\subsection{Detección de anomalías}

El módulo de detección de anomalías analiza la información almacenada con el objetivo de identificar comportamientos que se alejan de los patrones considerados normales.

Para cada anomalía detectada se calcula un nivel de riesgo que posteriormente es utilizado por el dashboard para clasificar y representar los resultados obtenidos.

La detección implementada permite identificar situaciones potencialmente peligrosas relacionadas con accesos no autorizados, movimientos en zonas protegidas, condiciones ambientales fuera de rango y fallos de dispositivos críticos.

La Figura D.6 muestra un ejemplo de las anomalías detectadas por el sistema durante las pruebas realizadas.

\imagen{anomalias}
{Resultado de la detección de anomalías clasificadas según su severidad y nivel de riesgo.}

\subsection{API REST}

La API REST desarrollada mediante FastAPI proporciona acceso a la información almacenada por el sistema. Los diferentes endpoints permiten consultar eventos, alertas, anomalías y parámetros de configuración.

La documentación interactiva generada automáticamente mediante Swagger facilita tanto las tareas de desarrollo como la validación de los servicios implementados.

La Figura D.7 muestra la documentación interactiva generada automáticamente por FastAPI, desde la cual es posible consultar y probar los diferentes servicios disponibles.

\imagen{swagger_api}
{Documentación interactiva de la API REST generada mediante Swagger.}

La Figura D.8 muestra un ejemplo de consulta realizada sobre uno de los endpoints implementados para la recuperación de alertas generadas por el sistema.

\imagen{swagger_alerts}
{Consulta realizada mediante el endpoint \textit{/alerts}.}

\subsection{Dashboard de monitorización}

La visualización de la información se realiza mediante una aplicación web desarrollada con Streamlit. El dashboard proporciona acceso a eventos, alertas, anomalías, sensores y parámetros de configuración del sistema.

Su diseño permite supervisar de forma centralizada el estado general de la instalación monitorizada y facilita la detección rápida de incidencias.

La Figura D.9 muestra la vista principal del dashboard desarrollado para el proyecto, donde se presentan indicadores generales del estado del sistema y estadísticas de funcionamiento.

\imagen{dashboard_principal}
{Vista principal del dashboard de monitorización.}

La Figura D.10 muestra una de las secciones destinadas a la consulta de alertas generadas por el sistema durante su funcionamiento.

\imagen{dashboard_2alertas}
{Visualización de alertas generadas por el sistema.}


\section{Compilación, instalación y ejecución del proyecto}

Una vez instaladas las dependencias necesarias, el sistema puede ejecutarse de forma manual o mediante el script de arranque desarrollado para el proyecto.

La ejecución manual de los diferentes componentes se realiza mediante los siguientes comandos:

\begin{verbatim}
	# Ejecución del núcleo principal del SIEM
	python main.py
	
	# Ejecución de la API REST
	python -m uvicorn api.server:app --reload
	
	# Ejecución del dashboard web
	python -m streamlit run dashboard/app.py
\end{verbatim}

Con el objetivo de simplificar la puesta en marcha del sistema se desarrolló el script \textit{start.bat}, que permite iniciar automáticamente el núcleo principal del SIEM, la API REST y el dashboard web desde una única ejecución.

La Figura D.11 muestra la ejecución del script de arranque desarrollado para el proyecto, encargado de iniciar automáticamente los principales componentes del sistema.

\imagen{start_bat}
{Ejecución del script de arranque \textit{start.bat}.}

Este mecanismo facilita el despliegue y las pruebas del sistema, evitando la necesidad de ejecutar cada componente de forma independiente y tener que tener los comandos de ejecución a mano, lo que lo simplifica significativamente su ejecución.

\section{Pruebas del sistema}

Durante el desarrollo se realizaron diferentes pruebas funcionales con el objetivo de verificar el correcto funcionamiento de cada uno de los módulos implementados y validar la integración entre los distintos componentes del sistema.

Las principales pruebas realizadas fueron:

\begin{itemize}
	\item Verificación de la lectura y normalización de eventos desde archivos de log.
	\item Comprobación del almacenamiento correcto de eventos y alertas en la base de datos SQLite.
	\item Validación de las reglas de correlación implementadas.
	\item Comprobación de la detección de anomalías y cálculo del nivel de riesgo.
	\item Verificación de los servicios expuestos mediante la API REST.
	\item Validación de la visualización de información en el dashboard desarrollado con Streamlit.
	\item Pruebas de supervisión del estado de sensores y detección de inactividad.
	\item Validación del mecanismo de simulación de notificaciones SMS para alertas críticas.
\end{itemize}

Además de las pruebas individuales realizadas sobre cada módulo, también se llevaron a cabo pruebas de integración para verificar el flujo completo de funcionamiento del sistema, desde la generación de eventos hasta su visualización en el dashboard.

Los resultados obtenidos fueron satisfactorios, verificándose el correcto funcionamiento de todos los componentes principales del sistema y la correcta interacción entre ellos.

A continuación se muestran algunas de las evidencias obtenidas durante las pruebas realizadas sobre los diferentes módulos del sistema.

\subsection{Pruebas del módulo de ingesta}
\imagen{ingestor}
{Ejecución del módulo de ingesta durante las pruebas realizadas.}

\subsection{Pruebas de la API REST}

La Figura D.12 muestra la documentación generada automáticamente mediante Swagger, utilizada para verificar el correcto funcionamiento de los endpoints implementados.

\imagen{swagger_api}
{Documentación Swagger generada automáticamente por FastAPI.}

La Figura D.13 muestra una consulta realizada sobre el endpoint de eventos.

\imagen{swagger_events}
{Consulta de eventos realizada mediante Swagger.}

\subsection{Pruebas de almacenamiento en SQLite}

Durante las pruebas se verificó el correcto almacenamiento de eventos y alertas dentro de la base de datos SQLite.

\imagen{sqlite_events}
{Contenido de la tabla \textit{events} almacenada en SQLite.}

\imagen{sqlite_alerts}
{Contenido de la tabla \textit{alerts} almacenada en SQLite.}

\subsection{Pruebas de detección de anomalías}

Se verificó el correcto funcionamiento del módulo de detección de anomalías mediante consultas realizadas a través de la API REST.

\imagen{anomalias}
{Resultado de la detección de anomalías clasificadas según su severidad y nivel de riesgo.}

\subsection{Pruebas del dashboard}

También se validó la correcta visualización de la información a través del dashboard desarrollado con Streamlit.

\imagen{dashboard_principal}
{Vista principal del dashboard.}

\imagen{dashboard_2alertas}
{Visualización de alertas generadas por el sistema.}

\imagen{dashboard_2eventos}
{Visualización de eventos almacenados.}

\subsection{Análisis de calidad del código}

Se utilizó la herramienta Codacy para realizar un análisis estático del código fuente y detectar posibles incidencias relacionadas con estilo, documentación, complejidad y buenas prácticas.

\imagen{codacy}
{Resultado del análisis estático realizado mediante Codacy.}

Los resultados obtenidos durante las pruebas permitieron verificar el correcto funcionamiento de todos los módulos implementados y confirmar la integración adecuada entre los diferentes componentes que forman parte del sistema SIEM.
